\documentclass[10pt,a4paper]{amsart}
\usepackage[margin=19mm]{geometry}
\usepackage{amsmath,amssymb,mathtools}
\usepackage[scr=rsfs,cal=euler]{mathalfa}
\usepackage{xfrac}
\usepackage[unicode,hidelinks,bookmarks=false]{hyperref}
\newcommand{\ie}{i.e.\ }
\newcommand{\cf}{cf.\ }
\newcommand{\resp}{respectively\ }
\newcommand{\Subsection}{\subsection}
\providecommand{\nobreakdash}{\nobreak\hbox{-}}
\setlength{\emergencystretch}{2em}
\multlinegap0pt
\usepackage{bm}
\usepackage{bbm}
\usepackage{dsfont}
\usepackage{xspace}
\usepackage{cancel}
\usepackage{mathtools}
\usepackage{amsthm}
\newtheorem{theorem}{Theorem}
\newtheorem{proposition}[theorem]{Proposition}
\newcommand{\dist}{\text{dist}}
\title[Asymmetric Homomorphism Thresholds for Graphs of Large Odd G]{Asymmetric Homomorphism Thresholds for Graphs of Large Odd Girth}
\author{Romain Bourneuf, Raphael Steiner, St\'ephan Thomass\'e, Yuval Wigderson}
\date{Source: arXiv:2609.25390v1 (CC BY 4.0)}
\begin{document}
\maketitle

\noindent\emph{Source and licence.} Asymmetric Homomorphism Thresholds for Graphs of Large Odd Girth. Version arXiv:2609.25390v1 (CC BY 4.0), distributed under the Creative Commons Attribution 4.0 International licence, as recorded in the arXiv licence metadata for this exact version and shown on the abstract page https://arxiv.org/abs/2609.25390v1. Full rights notice: licenses/arxiv-2609.25390-CC-BY-4.0.txt.\par\smallskip

\noindent\textit{Editorial scope.} Counted unit: the Proposition whose source label is \texttt{prop:construction} in the article (source0.tex lines 330--336, source label \texttt{prop:construction}), delivered together with the article's own complete original proof of it (source0.tex lines 338--397, one \texttt{proof} environment of 60 lines). No mathematics is written, repaired or re-derived: statement and proof are the article's own text line for line. Required context retained uncounted: none. Every definition, display and label of those lines is retained unchanged. Omitted from this package and not redistributed: 1--329, 337--337, 398--685. Nothing inside a retained excerpt is omitted or altered: every retained body is delivered exactly as the article's lines are written, so no omission receipt of any kind accompanies this package. Citations: the retained excerpts carry no citation command at all, so no bibliography entry of the article is transcribed and no reference is redistributed. Reference adaptation: none was needed and none was made; every reference inside the delivered unit resolves inside it, so no unresolved reference or citation is produced. Printed slips kept literal: no printed word, symbol, hypothesis or number of the counted statement or of its proof was altered. Layout and class: the article's own class is \texttt{article} with packages including graphicx, amsmath, amssymb, amsthm, thmtools, xcolor, xspace, tcolorbox, mathtools, enumitem, thm-restate, babel, geometry, hyperref; the wrapper is the lane's pinned \texttt{amsart} editorial wrapper at 10pt on A4, which restates the theorem-like environments the retained text needs and adds the article's own macro definitions that the retained text needs (\texttt{\textbackslash dist}), with extra wrapper packages (bm, bbm, dsfont, xspace, cancel, mathtools). The delivered numbering is the unit's own. Assets: no retained span contains a figure environment, a graphics-inclusion command, a PDF-page inclusion, a table or a TikZ picture; no image file is copied into this package, the package declares no asset dependency and no image file is redistributed with it. Licence: version arXiv:2609.25390v1 of this article is distributed under the Creative Commons Attribution 4.0 International licence, as recorded in the arXiv licence metadata for this exact version and shown on the abstract page https://arxiv.org/abs/2609.25390v1. A full rights notice is delivered at licenses/arxiv-2609.25390-CC-BY-4.0.txt. Characters: every retained excerpt is pure ASCII, so the delivered engine, XeLaTeX with its default Latin Modern text font, reports no missing character.
\begin{proposition}\label{prop:construction}
    For every integer $N \geq 1$, for every $\varepsilon > 0$, and for every graph $H$ which contains an odd cycle, there exists a graph $H'$ with the following properties: \begin{itemize}
        \item $og(H') \geq og(H) + 4$,
        \item $\delta(H') \geq \left(\frac{1}{3|V(H)|} - \varepsilon\right)|V(H')|$,
        \item If $H'$ has a homomorphism to a graph $J$ on at most $N$ vertices then $H$ also has a homomorphism to $J$.
    \end{itemize}
\end{proposition}

\begin{proof}
    Let $d \in \mathbb{N}$ be large enough so that for every $\theta > 0$, the Borsuk graph $Bor(d, \theta)$ has chromatic number greater than $N^{|V(H)|}$.
    Define $\zeta = \varepsilon|V(H)|$ and let $\eta>0$ be chosen small enough so that $\frac{1}{2+\eta}\cdot \frac{1}{3/2 + \eta} \geq \frac{1}{3} - \zeta$.
    Let $\theta \in (0, \pi)$ be small enough so that $\chi_f(Bor(d, \theta)) \leq 2 + \eta$ and let $\iota < \frac{\theta}{og(H)+2}$.
    By the \emph{de Bruijn--Erd\H{o}s theorem}, there exists a finite set $X \subseteq \mathbb{S}^d$ such that $\chi(Bor(d, \iota)[X]) = \chi(Bor(d, \iota)) > N^{|V(H)|}$.
    Let $F \coloneqq Bor(d, \theta)[X]$, and observe that $\chi_f(F) \leq \chi_f(Bor(d, \theta)) \leq 2 + \eta$.

    Let $\mathcal{I}$ be the set of all independent sets of $F$.
    Since $\chi_f(F) \leq 2 + \eta$, there exists a weight function $w : \mathcal{I} \to \mathbb{R}^+$ such that $\sum_{v \in I \in \mathcal{I}}w(I) \geq 1$ for every $v \in X$ and $w(\mathcal{I}) \coloneqq \sum_{I \in \mathcal{I}}w(I) \leq 2+\eta$.
    Since $\chi_f(F)$ is the solution of a (finite) linear program with integral coefficients, we may assume without loss of generality that $w : \mathcal{I} \to \mathbb{Q}^+$.
    Thus, after rescaling, there exists a function $w' : \mathcal{I} \to \mathbb{N}$ such that $\sum_{v \in I \in \mathcal{I}}w'(I) \geq \frac{1}{2+\eta} \cdot w'(\mathcal{I})$ for every $v \in X$.
    By considering the (bipartite) incidence graph between the vertices $X$ and the independent sets $\mathcal{I}$, and replacing each vertex corresponding to an independent set $I \in \mathcal{I}$ by $w'(I)$ copies of itself, it follows that there exists a set $Y$ and a bipartite graph $B$ with bipartition $(X, Y)$ such that $|N_B(v)| \geq \frac{1}{2+\eta} \cdot |Y|$ for every $v \in X$.
    Note that $Y$ can be chosen to have even size and arbitrarily large, by replacing each element of $Y$ by a large (but fixed) number of copies of itself.
    In particular, we can assume that $|Y|$ is even and that $|X| \leq \eta|Y|$.
    By adding a set $Z$ of size $|Y|/2$ that is complete to $Y$, we obtain a bipartite graph $G$ with bipartition $(X \cup Z, Y)$.
    Note that $|V(G)| = |X| + |Z| + |Y| = |X| + \frac{3}{2}|Y| \leq \left(\frac{3}{2} + \eta\right)|Y|$.
    Then, every $v \in X \cup Z$ satisfies $|N_{G}(v)| \geq \frac{1}{2+\eta} \cdot |Y| \geq \frac{1}{2 + \eta} \cdot \frac{1}{3/2 + \eta}|V(G)| \geq  \left(\frac{1}{3} - \zeta\right)|V(G)|$ by definition of $\eta$.
    Also, every $v \in Y$ satisfies $|N_{G}(v)| \geq |Z| = \frac{|Y|}{2} \geq \frac{1}{2} \cdot \frac{|V(G)|}{3/2 + \eta} \geq \frac{1}{2 + \eta} \cdot \frac{1}{3/2 + \eta}|V(G)| \geq \left(\frac{1}{3} - \zeta\right)|V(G)|$ again by definition of $\eta$.
    Thus, $\delta(G) \geq \left(\frac{1}{3} - \zeta\right)|V(G)|$.

    For every $h \in V(H)$, let $G_h$ be a copy of $G$ such that all copies $G_h$ are vertex-disjoint.
    For every $h \in V(H)$, denote the vertex set of $G_h$ by $X_h \cup Y_h \cup Z_h$, and for every $x \in X$ let $x_h$ denote the copy of $x$ in $X_h$.
    Then, let $H'$ be the graph obtained from the disjoint union of all $G_h$ by adding edges between the sets $X_h$ as in the tensor (or categorical) product $Bor(d, \iota) \times H$.
    Formally, connect $x_h$ and $x'_{h'}$ for $x, x' \in X$ and $h, h' \in V(H)$ if and only if $xx' \in E(Bor(d, \iota))$ (i.e. if $\dist(x, x') \geq \pi - \iota$) and $hh' \in E(H)$.

    We now show that $H'$ satisfies the desired properties.
    By contradiction, suppose that $og(H') < og(H) + 4$, so $og(H') \leq og(H) + 2$.
    Since we can always extend the length of a closed walk by any even number by moving forth and back along some edge, it follows that there exists a closed walk $W'$ of length \emph{exactly} $\ell \coloneqq og(H) + 2$ in $H'$.
    Let $H^+$ be the graph obtained from $H$ by adding for every $h \in V(H)$ a disjoint edge $e_h \coloneqq a_hb_h$ which we connect to $h$ via another edge $ha_h$. 
    By construction, the map $\phi : V(H') \to V(H^+)$ that maps each set $X_h$ to $h$, each set $Y_h$ to $a_h$ and each set $Z_h$ to $b_h$, is a homomorphism from $H'$ to $H^+$.
    Thus, $\phi$ maps $W'$ to a closed walk $W^+$ of length $\ell$ in $H^+$.
    Suppose first that $W^+$ only visits vertices of $H$ in $H^+$.
    Then, we can write $W' = x^0_{h_0}x^1_{h_1}\ldots x^{\ell-1}_{h_{\ell-1}}x^0_{h_0}$ where $h_i \in V(H)$ and $x_i \in X$ for every $i \in \{0, \ldots, \ell\}$.
    A simple proof by induction shows that $\dist(x^0, x^j) \leq j\iota$ for every even $j \in \{0, \ldots, \ell-1\}$.
    Since $\ell-1$ is even, we have $\dist(x^0, x^{\ell-1}) \leq (\ell-1)\iota$.
    However, since $x^{\ell-1}_{h_{\ell-1}}x^0_{h_0}$ is an edge of $H'$, we have $\dist(x^{0}, x^{\ell-1}) \geq \pi-\iota$, so $\pi-\iota \leq (\ell-1)\iota$ so $\iota \geq \frac{\pi}{\ell} = \frac{\pi}{og(H)+2}$, a contradiction since $\iota < \frac{\theta}{og(H)+2} < \frac{\pi}{og(H)+2}$ by definition.
    Thus, $W^+$ does not only visit vertices from $V(H)$ in $H^+$.
    Since $W^+$ has length $og(H)+2$ and every odd cycle of $H^+$ lives in $V(H)$, so has length at least $og(H)$, it follows that there exists $h \in V(H)$ so that $W^+$ visits $h$, then $a_h$, then $h$ again, and then forms a cycle of length $og(H)$ that only visits vertices of $H$ in $H^+$.
    Let $x^0_h$ and $x^2_h$ be the two vertices visited by $W'$ that belong to $X_h$.
    Then, $x^0_h$ and $x^2_h$ have a common neighbor in $Y_h$, and so there is an independent set in $F = Bor(d, \theta)[X]$ containing both $x^0$ and $x^2$. Hence, $x^0$ and $x^2$ are not adjacent in $F = Bor(d, \theta)[X]$, so $\dist(x^0, x^2) \leq \pi - \theta$.
    However, as before, a simple proof by induction shows that $\dist(x^0, x^2) \geq \pi - og(H)\iota$.
    This implies $\pi-og(H)\iota \leq \pi - \theta$, so $\iota \geq \frac{\theta}{og(H)}$, again a contradiction.
    Thus, $og(H') \geq og(H) + 4$.
    

    By definition, we have $|V(H')| = |V(H)| \cdot |V(G)|$.
    Let $v \in V(H')$ and let $h \in V(H)$ such that $v \in V(G_h)$.
    Then, \[d_{H'}(v) \geq d_{G_h}(v) \geq \left(\frac{1}{3} - \zeta\right)|V(G)| = \left(\frac{1}{3|V(H)|} - \frac{\zeta}{|V(H)|}\right)|V(H')| = \left(\frac{1}{3|V(H)|} - \varepsilon\right)|V(H')|.\]

    Let $J$ be a graph on at most $N$ vertices such that there exists a homomorphism $\phi : V(H') \to V(J)$ from $H'$ to $J$.
    Without loss of generality, let us assume that $V(J) = [N]$.
    Write $V(H) = \{h_1, \ldots, h_{|V(H)|}\}$.
    For each $x \in X$, let $c(x) \coloneqq (\phi(x_{h_1}), \ldots, \phi(x_{h_{|V(H)|}})) \in [N]^{|V(H)|}$.
    Since we have $\chi(Bor(d, \iota)[X]) > N^{|V(H)|}$, it follows that $c$ is not a proper coloring of $Bor(d, \iota)[X]$, so there exist $x, y \in X$ such that $xy \in E(Bor(d, \iota)[X])$ and $c(x) = c(y)$.
    We show that the map $\psi : V(H) \to [N], h \mapsto \phi(x_h)$ is a homomorphism from $H$ to $J$.
    Let $h, h' \in V(H)$ such that $hh' \in E(H)$.
    Then, $x_hy_{h'} \in E(H')$ so $\phi(x_h)\phi(y_{h'})$ is an edge of $J$.
    By definition, we have $\psi(h) = \phi(x_h)$ and $\psi(h') = \phi(x_{h'}) = \phi(y_{h'})$ since $c(x) = c(y)$.
    Thus, $\psi(h)\psi(h') \in E(J)$, so $\psi$ is indeed a homomorphism.
\end{proof}

\end{document}
